\begin{abstract}
Temporal knowledge graph forecasting asks which object completes a query $(s, r, ?, t)$ given every fact before $t$. The strongest published models answer it by reading the history of the query pair $(s, r)$: a copy mechanism over its past objects, or a graph built from its past facts. We show that this is the wrong unit of evidence. The unit that carries the signal is the \emph{dyad}, the ordered entity pair $(s, o)$ together with the typed, timestamped stream of every event between them under any relation and in either direction. On ICEWS18, 23\% of test queries have an answer that never occurred with $(s, r)$ yet had interacted with $s$ before; every copy mechanism is silent on them by construction. A counter that merely reads the dyad, with no trained parameter, already matches DaeMon on that benchmark.

We present \model{}, which treats a query as a draw from a superposition of four marked point processes whose intensities are added, never gated: a structural intensity from snapshot evolution, a popularity intensity over all entities, a dyadic intensity in which a small Transformer encodes the dyad's event stream, and a path intensity that forwards each dyadic state along the candidate's own most recent typed edges. Candidates of one query attend to one another before they are scored, so an intensity can depend on what the competing candidates' streams look like.

Over three seeds and five benchmarks, \model{} sets the best published MRR and \hone{} on YAGO (91.82 / 90.57) and WIKI (83.41 / 80.44), the best \hone{} on GDELT (18.16), and on ICEWS18 and ICEWS14s exceeds every model whose code we could run under one protocol. A stratified evaluation, defined from the data alone, shows where each intensity acts: the dyadic stream and its relation types matter on event data and not on yearly data, exactly where the dyad stratum exists. Code, result files and the diagnostic are released.
\end{abstract}
